% GENERATED FILE. Edit canonical lessons, then run npm run generate:books.
% canonical-source-hash: fnv1a64:0c6bdd27b96e0f2d
% canonical-lessons: PT-C333-batedeira, PT-C333-peneirar, PT-C333-amassar, PT-C333-untar, PT-C333-polvilhar

\chapter{Um bolo caseiro}
\label{ch:pt-um-bolo-caseiro}
\hlchaptermodality{333}

\begin{chapteropening}
\textbf{By the end of this chapter:} \emph{I can describe baking a cake: the mixer, sifting, kneading, greasing the tin and dusting.}

Complete the last lesson of chapter 333: I can describe baking a cake: the mixer, sifting, kneading, greasing the tin and dusting.
\end{chapteropening}

\section[a batedeira]{a batedeira — the electric mixer}

\label{lesson:PT-C333-batedeira}



\begin{quote}
Before the new one: say the Portuguese for to grill, then the Portuguese for to boil.
\end{quote}

\subsection*{You'll want to know: a batedeira}
\textbf{a batedeira} — \textquotedblleft{}the electric mixer\textquotedblright{}.

\begin{quote}
A minha mãe bate as claras na batedeira.
\end{quote}

\begin{quote}
\emph{My mother beats the egg whites in the mixer.}
\end{quote}

Built on \emph{bater}, ‘to beat’.

\begin{cousinweb}
\emph{batedeira} is built on \emph{bater} \textquotedblleft{}to beat\textquotedblright{}, which comes from Latin \emph{battuere} \textquotedblleft{}to beat\textquotedblright{} (Lewis \& Short), with the ending \emph{-deira} for a tool.
\end{cousinweb}

\begin{grammarlens}[title={What you've built}]
One more word for telling, explaining and describing at B1.
\end{grammarlens}

\subsection*{Your turn}
\begin{itemize}
\raggedright
  \item \emph{Say it:} \emph{a batedeira}
  \item \emph{Say it:} \emph{a batedeira}, once more
  \item \emph{Say it:} a sentence of your own with \emph{a batedeira}
  \item \emph{Recall:} say the Portuguese for to grill, then the Portuguese for to boil, then say \emph{a batedeira} again
\end{itemize}

\begin{tcolorbox}[breakable,colback=teal!4,colframe=teal!35!black,title={Before you move on}]
What does a batedeira mean? (\textquotedblleft{}The electric mixer\textquotedblright{}.) Say it once more.
\end{tcolorbox}
\section[peneirar]{peneirar — to sift}

\label{lesson:PT-C333-peneirar}



\begin{quote}
Before the new one: say the Portuguese for to boil, then the Portuguese for the electric mixer.
\end{quote}

\subsection*{You'll want to know: peneirar}
\textbf{peneirar} — \textquotedblleft{}to sift\textquotedblright{}.

\begin{quote}
A avó peneira a farinha duas vezes.
\end{quote}

\begin{quote}
\emph{Grandmother sifts the flour twice.}
\end{quote}

Built on \emph{a peneira}, ‘the sieve’.

\begin{cousinweb}
\emph{peneirar} is built on \emph{peneira} \textquotedblleft{}a sieve\textquotedblright{}, which is inherited from Latin \emph{pānāria} (EtymDB), akin to \emph{pānārium} \textquotedblleft{}a bread basket\textquotedblright{}, from \emph{pānis} \textquotedblleft{}bread\textquotedblright{} (Lewis \& Short).
\end{cousinweb}

\begin{grammarlens}[title={What you've built}]
One more word for telling, explaining and describing at B1.
\end{grammarlens}

\subsection*{Your turn}
\begin{itemize}
\raggedright
  \item \emph{Say it:} \emph{peneirar}
  \item \emph{Say it:} \emph{peneirar}, once more
  \item \emph{Say it:} a sentence of your own with \emph{peneirar}
  \item \emph{Recall:} say the Portuguese for to boil, then the Portuguese for the electric mixer, then say \emph{peneirar} again
\end{itemize}

\begin{tcolorbox}[breakable,colback=teal!4,colframe=teal!35!black,title={Before you move on}]
What does peneirar mean? (\textquotedblleft{}To sift\textquotedblright{}.) Say it once more.
\end{tcolorbox}
\section[amassar]{amassar — to knead; to crush, to dent}

\label{lesson:PT-C333-amassar}



\begin{quote}
Before the new one: say the Portuguese for the electric mixer, then the Portuguese for to sift.
\end{quote}

\subsection*{You'll want to know: amassar}
\textbf{amassar} — \textquotedblleft{}to knead; to crush, to dent\textquotedblright{}.

\begin{quote}
O padeiro amassa o pão de madrugada.
\end{quote}

\begin{quote}
\emph{The baker kneads the bread before dawn.}
\end{quote}

Also to dent a car: \emph{amassar o carro}.

\begin{cousinweb}
\emph{amassar} is built from \emph{a-} and \emph{massa} \textquotedblleft{}dough\textquotedblright{}, from Latin \emph{massa} \textquotedblleft{}a lump, a mass\textquotedblright{} (Lewis \& Short).
\end{cousinweb}

\begin{grammarlens}[title={What you've built}]
One more word for telling, explaining and describing at B1.
\end{grammarlens}

\subsection*{Your turn}
\begin{itemize}
\raggedright
  \item \emph{Say it:} \emph{amassar}
  \item \emph{Say it:} \emph{amassar}, once more
  \item \emph{Say it:} a sentence of your own with \emph{amassar}
  \item \emph{Recall:} say the Portuguese for the electric mixer, then the Portuguese for to sift, then say \emph{amassar} again
\end{itemize}

\begin{tcolorbox}[breakable,colback=teal!4,colframe=teal!35!black,title={Before you move on}]
What does amassar mean? (\textquotedblleft{}To knead; to crush, to dent\textquotedblright{}.) Say it once more.
\end{tcolorbox}
\section[untar]{untar — to grease (a tin), to rub with fat}

\label{lesson:PT-C333-untar}



\begin{quote}
Before the new one: say the Portuguese for to sift, then the Portuguese for to knead.
\end{quote}

\subsection*{You'll want to know: untar}
\textbf{untar} — \textquotedblleft{}to grease (a tin), to rub with fat\textquotedblright{}.

\begin{quote}
Antes de pôr a massa, ela unta a forma com manteiga.
\end{quote}

\begin{quote}
\emph{Before putting the mixture in, she greases the tin with butter.}
\end{quote}

\emph{Untar a forma} — to grease the tin.

\begin{cousinweb}
\emph{untar} comes down from Old Portuguese (EtymDB) and goes back to Latin \emph{ūnctus} \textquotedblleft{}anointed\textquotedblright{}, from \emph{unguere} \textquotedblleft{}to smear with oil\textquotedblright{} (Lewis \& Short).
\end{cousinweb}

\begin{grammarlens}[title={What you've built}]
One more word for telling, explaining and describing at B1.
\end{grammarlens}

\subsection*{Your turn}
\begin{itemize}
\raggedright
  \item \emph{Say it:} \emph{untar}
  \item \emph{Say it:} \emph{untar}, once more
  \item \emph{Say it:} a sentence of your own with \emph{untar}
  \item \emph{Recall:} say the Portuguese for to sift, then the Portuguese for to knead, then say \emph{untar} again
\end{itemize}

\begin{tcolorbox}[breakable,colback=teal!4,colframe=teal!35!black,title={Before you move on}]
What does untar mean? (\textquotedblleft{}To grease (a tin), to rub with fat\textquotedblright{}.) Say it once more.
\end{tcolorbox}
\section[polvilhar]{polvilhar — to sprinkle, to dust}

\label{lesson:PT-C333-polvilhar}



\begin{quote}
Before the new one: say the Portuguese for to knead, then the Portuguese for to grease.
\end{quote}

\subsection*{You'll want to know: polvilhar}
\textbf{polvilhar} — \textquotedblleft{}to sprinkle, to dust\textquotedblright{}.

\begin{quote}
No fim, ele polvilha o bolo com açúcar.
\end{quote}

\begin{quote}
\emph{At the end he dusts the cake with sugar.}
\end{quote}

Of a powder: sugar, flour, cinnamon.

\begin{cousinweb}
\emph{polvilhar} is built on \emph{polvilho} \textquotedblleft{}fine powder\textquotedblright{}, a word built on \emph{pó} \textquotedblleft{}dust\textquotedblright{}, from Latin \emph{pulvis} \textquotedblleft{}dust\textquotedblright{} (Lewis \& Short).
\end{cousinweb}

\begin{grammarlens}[title={What you've built}]
One more word for telling, explaining and describing at B1.
\end{grammarlens}

\subsection*{Your turn}
\begin{itemize}
\raggedright
  \item \emph{Say it:} \emph{polvilhar}
  \item \emph{Say it:} \emph{polvilhar}, once more
  \item \emph{Say it:} a sentence of your own with \emph{polvilhar}
  \item \emph{Recall:} say the Portuguese for to knead, then the Portuguese for to grease, then say \emph{polvilhar} again
\end{itemize}

\begin{tcolorbox}[breakable,colback=teal!4,colframe=teal!35!black,title={Before you move on}]
What does polvilhar mean? (\textquotedblleft{}To sprinkle, to dust\textquotedblright{}.) Say it once more.
\end{tcolorbox}
